\chapter{Результаты вычислений и программная реализация}

Программный комплекс для решения поставленных задач разработан на языке программирования C++ и скомпонован в файле \texttt{main.cpp}.

\section{Структура программы и ввод данных}

Для универсальности работы с системами произвольной размерности используются стандартные контейнеры библиотеки STL: динамические массивы \texttt{std::vector<double>} (тип \texttt{vd}) и матрицы \texttt{std::vector<std::vector<double>>} (тип \texttt{vvd}).

Считывание коэффициентов системы осуществляется функцией \texttt{readFile} (листинг~\ref{lst:code-read}). Функция перенаправляет стандартный поток ввода на целевой файл и выполняет построчный разбор с помощью строкового потока \texttt{std::stringstream}, что позволяет автоматически определять размерность системы без жесткого задания констант.

\begin{listing}[H]
    \cppfile[firstline=18,lastline=36,firstnumber=18,autogobble]{../main.cpp}
    \caption{Функция readFile для универсального считывания матрицы}
    \label{lst:code-read}
\end{listing}

\section{Реализация метода Гаусса с выбором главных элементов}

\subsection{Программные модули прямого хода}

Поиск главного элемента по всей активной подматрице реализован в функции \texttt{findMajor} (листинг~\ref{lst:code-findmajor}). Функция сканирует все элементы подматрицы и находит элемент с максимальной абсолютной величиной. В строке~56 выполняется проверка превосходства порога машинного нуля $\varepsilon_{\text{mach}} = 10^{-9}$, что предотвращает деление на ноль при вырожденности системы.

\begin{listing}[H]
    \cppfile[firstline=47,lastline=65,firstnumber=47,autogobble]{../main.cpp}
    \caption{Функция findMajor для поиска максимального по модулю элемента}
    \label{lst:code-findmajor}
\end{listing}

Исключение переменной из остальных строк системы выполняется функцией \texttt{sumRows} (листинг~\ref{lst:code-sumrows}). В строке~75 вычисляется масштабирующий множитель $m_i = -(a_{iq} / a_{pq})$, после чего к каждой строке $i \neq p$ прибавляется главная строка, умноженная на $m_i$.

\begin{listing}[H]
    \cppfile[firstline=67,lastline=79,firstnumber=67,autogobble]{../main.cpp}
    \caption{Функция sumRows для исключения главного столбца из строк}
    \label{lst:code-sumrows}
\end{listing}

Для сокращения размерности подматрицы на каждом шаге алгоритма используется вспомогательная функция \texttt{deleteMajor} (листинг~\ref{lst:code-deletemajor}), которая заполняет нулями главную строку и главный столбец, формируя матрицу для следующего рекурсивного вызова.

\begin{listing}[H]
    \cppfile[firstline=81,lastline=98,firstnumber=81,autogobble]{../main.cpp}
    \caption{Функция deleteMajor для обнуления главной строки и столбца}
    \label{lst:code-deletemajor}
\end{listing}

Основная логика решения СЛАУ методом Гаусса сосредоточена в рекурсивной функции \texttt{task1} (листинг~\ref{lst:code-task1}). Если функция \texttt{findMajor} не находит ненулевого элемента (строки~106--111), выполняется анализ свободных членов: при наличии ненулевого остатка фиксируется несовместность (возвращается пустой вектор), а при нулевом остатке возвращается признак бесконечного множества решений. Если система совместна, в строках~127--131 выполняется обратный ход Гаусса с вычислением неизвестной $x_q$.

\begin{listing}[H]
    \cppfile[firstline=100,lastline=134,firstnumber=100,autogobble]{../main.cpp}
    \caption{Рекурсивная функция task1 метода Гаусса с обратным ходом}
    \label{lst:code-task1}
\end{listing}

\subsection{Пошаговый протокол выполнения метода Гаусса}

В результате работы функции \texttt{task1} над системой варианта №~11 зафиксированы следующие промежуточные этапы:

\textbf{Шаг 1.} Поиск максимума среди всех коэффициентов исходной матрицы:
\begin{equation*}
    \max_{i, j} |a_{ij}| = |a_{23}| = 6{,}7.
\end{equation*}
Главный элемент расположен во 2-й строке и 3-м столбце (соответствует переменной $x_3$). Вычисляются множители для строк 1 и 3:
\begin{equation*}
    m_1 = -\frac{-1{,}5}{6{,}7} \approx 0{,}223881, \qquad m_3 = -\frac{-1{,}4}{6{,}7} \approx 0{,}208955.
\end{equation*}
После сложения строк получаем преобразованную матрицу:
\begin{equation*}
    \begin{pmatrix}
        3{,}70746 & 0{,}273134 & 0 & | & 4{,}08209 \\
        4{,}5 & -2{,}8 & 6{,}7 & | & 2{,}6 \\
        6{,}04030 & 3{,}11493 & 0 & | & 0{,}403284
    \end{pmatrix}.
\end{equation*}
Строка 2 и столбец 3 исключаются из дальнейшего рассмотрения прямого хода.

\textbf{Шаг 2.} Поиск главного элемента среди оставшихся активных коэффициентов (строки 1 и 3, столбцы 1 и 2):
\begin{equation*}
    \max |a_{ij}| = |a_{31}| = 6{,}04030.
\end{equation*}
Главный элемент расположен в 3-й строке и 1-м столбце (переменная $x_1$). Множитель для строки 1:
\begin{equation*}
    m_1 = -\frac{3{,}70746}{6{,}04030} \approx -0{,}613788.
\end{equation*}
После сложения первой строки с третьей, умноженной на $m_1$, получаем:
\begin{equation*}
    \begin{pmatrix}
        0 & -1{,}63877 & 0 & | & 3{,}83456 \\
        0 & 0 & 0 & | & 0 \\
        6{,}04030 & 3{,}11493 & 0 & | & 0{,}403284
    \end{pmatrix}.
\end{equation*}
Строка 3 и столбец 1 исключаются.

\textbf{Шаг 3.} В активной подматрице остается единственный ненулевой элемент:
\begin{equation*}
    a_{12} = -1{,}63877 \quad (\text{переменная } x_2), \qquad b_1 = 3{,}83456.
\end{equation*}

\textbf{Обратный ход}:
\begin{enumerate}
    \item Из последнего уравнения находим $x_2$:
    \begin{equation*}
        x_2 = \frac{3{,}83456}{-1{,}63877} \approx -2{,}33990;
    \end{equation*}
    \item Из главной строки шага 2 (строка 3) находим $x_1$:
    \begin{equation*}
        x_1 = \frac{0{,}403284 - 3{,}11493 \cdot (-2{,}33990)}{6{,}04030} = \frac{7{,}69190}{6{,}04030} \approx 1{,}27343;
    \end{equation*}
    \item Из главной строки шага 1 (строка 2) находим $x_3$:
    \begin{equation*}
        x_3 = \frac{2{,}6 - 4{,}5 \cdot 1{,}27343 - (-2{,}8) \cdot (-2{,}33990)}{6{,}7} = \frac{-9{,}68215}{6{,}7} \approx -1{,}44510.
    \end{equation*}
\end{enumerate}

Итоговый вектор решения задачи 1 методом Гаусса:
\begin{equation}
    x_{\text{Gauss}} = \begin{pmatrix} 1{,}27343 \\ -2{,}33990 \\ -1{,}44510 \end{pmatrix}.
    \label{eq:sol_gauss}
\end{equation}

\section{Реализация метода простой итерации}

\subsection{Программные модули обеспечения сходимости}

Проверка выполнения условия строгого диагонального преобладания реализована в функции \texttt{isDominant} (листинг~\ref{lst:code-isdominant}).

\begin{listing}[H]
    \cppfile[firstline=158,lastline=167,firstnumber=158,autogobble]{../main.cpp}
    \caption{Функция isDominant для проверки диагонального преобладания}
    \label{lst:code-isdominant}
\end{listing}

Если диагональное преобладание отсутствует, вызывается функция \texttt{makeDominant} (листинг~\ref{lst:code-makedomiant}). Функция выполняет полный перебор целочисленных весов $w_r \in \{-1, 0, 1\}$ линейных комбинаций уравнений, находя такую эквивалентную систему, в которой диагональные коэффициенты строго превосходят сумму модулей недиагональных элементов.

\begin{listing}[H]
    \cppfile[firstline=169,lastline=199,firstnumber=169,autogobble]{../main.cpp}
    \caption{Функция makeDominant для синтеза диагонально доминирующей матрицы}
    \label{lst:code-makedomiant}
\end{listing}

Для проверки достаточного условия сходимости метода итераций реализованы функции вычисления норм матрицы \texttt{normOct} (октаэдрическая норма, листинг~\ref{lst:code-normoct}) и \texttt{normCub} (кубическая норма, листинг~\ref{lst:code-normcub}).

\begin{listing}[H]
    \cppfile[firstline=136,lastline=145,firstnumber=136,autogobble]{../main.cpp}
    \caption{Функция normOct для вычисления октаэдрической нормы матрицы}
    \label{lst:code-normoct}
\end{listing}

\begin{listing}[H]
    \cppfile[firstline=147,lastline=156,firstnumber=147,autogobble]{../main.cpp}
    \caption{Функция normCub для вычисления кубической нормы матрицы}
    \label{lst:code-normcub}
\end{listing}

Реализация итерационного процесса приведена в функции \texttt{task2}. Первая часть функции отвечает за приведение матрицы к виду со строгим диагональным преобладанием, построение матрицы перехода $B$ и вектора $c$, а также проверку условий сходимости по вычисленным нормам (листинг~\ref{lst:code-task2-prep}).

\begin{listing}[H]
    \cppfile[firstline=201,lastline=234,firstnumber=201,autogobble]{../main.cpp}
    \caption{Подготовка матрицы B и проверка сходимости в функции task2}
    \label{lst:code-task2-prep}
\end{listing}

Вторая часть функции \texttt{task2} реализует непосредственно итерационный цикл (листинг~\ref{lst:code-task2-loop}). Согласно требованиям к заданию 2, для хранения неизвестных используются строго два одномерных массива: \texttt{x\_prev} (вектор значений на предыдущей итерации) и \texttt{x\_curr} (вектор значений на текущей итерации). Итерационный цикл завершается при достижении условия $\max_i |x_i^{(k)} - x_i^{(k-1)}| < 10^{-5}$.

\begin{listing}[H]
    \cppfile[firstline=235,lastline=266,firstnumber=235,autogobble]{../main.cpp}
    \caption{Итерационный цикл с двумя одномерными массивами в функции task2}
    \label{lst:code-task2-loop}
\end{listing}

\subsection{Пошаговый протокол выполнения метода простой итерации}

После применения функции \texttt{makeDominant} получена эквивалентная система с диагональным преобладанием:
\begin{equation}
    \widetilde{A}_{\text{dom}} = \begin{pmatrix}
        -12{,}3 & -1{,}8 & -3{,}8 & | & -5{,}96 \\
        2{,}4 & 2{,}8 & 0{,}1 & | & -3{,}64 \\
        1{,}8 & -3{,}7 & 8{,}2 & | & -0{,}90
    \end{pmatrix}.
    \label{eq:matrix_dom}
\end{equation}

Строки матрицы~\eqref{eq:matrix_dom} получены следующими линейными комбинациями исходных уравнений $R_1, R_2, R_3$:
\begin{enumerate}
    \item новая строка 1: $-R_1 - R_2 - R_3$, где $|-12{,}3| = 12{,}3 > |-1{,}8| + |-3{,}8| = 5{,}6$;
    \item новая строка 2: $-R_1 + R_3$, где $|2{,}8| = 2{,}8 > |2{,}4| + |0{,}1| = 2{,}5$;
    \item новая строка 3: $-R_1 + R_2$, где $|8{,}2| = 8{,}2 > |1{,}8| + |-3{,}7| = 5{,}5$.
\end{enumerate}

Выражая диагональные неизвестные, получаем каноническую систему $x = Bx + c$:
\begin{equation}
    B = \begin{pmatrix}
        0 & -0{,}146341 & -0{,}308943 \\
        -0{,}857143 & 0 & -0{,}035714 \\
        -0{,}219512 & 0{,}451220 & 0
    \end{pmatrix},
    \qquad
    c = \begin{pmatrix}
        0{,}484553 \\
        -1{,}300000 \\
        -0{,}109756
    \end{pmatrix}.
\end{equation}

Вычисление подчиненных норм матрицы перехода $B$:
\begin{enumerate}
    \item октаэдрическая норма (максимум по столбцам):
    \begin{multline*}
        \|B\|_1 = \max\big(0 + 0{,}857143 + 0{,}219512;\, \\
        0{,}146341 + 0{,}451220;\, 0{,}308943 + 0{,}035714\big) = 1{,}07666;
    \end{multline*}
    \item кубическая норма (максимум по строкам):
    \begin{multline*}
        \|B\|_\infty = \max\big(0{,}146341 + 0{,}308943;\, \\
        0{,}857143 + 0{,}035714;\, 0{,}219512 + 0{,}451220\big) = 0{,}892857.
    \end{multline*}
\end{enumerate}

Так как кубическая норма матрицы $\|B\|_\infty = 0{,}892857 < 1$, \textbf{достаточное условие сходимости метода выполнено}.

Динамика приближений на каждом шаге итерационного процесса приведена в таблице~\ref{tab:iterations}.

\clearpage
\begin{table}[H]
    \caption{Пошаговый протокол выполнения метода простой итерации}
    \label{tab:iterations}
    \centering
    \scriptsize
    \renewcommand{\arraystretch}{0.85}
    \begin{tabular}{|c|c|c|c|c|}
        \hline
        \textbf{Итерация $k$} & $x_1^{(k)}$ & $x_2^{(k)}$ & $x_3^{(k)}$ & $\|x^{(k)} - x^{(k-1)}\|_\infty$ \\
        \hline
        $0$ (вектор $c$) & $0{,}484553$ & $-1{,}300000$ & $-0{,}109756$ & --- \\
        $1$ & $0{,}708705$ & $-1{,}711410$ & $-0{,}802707$ & $0{,}692951$ \\
        $2$ & $0{,}982994$ & $-1{,}878790$ & $-1{,}037550$ & $0{,}274289$ \\
        $3$ & $1{,}080040$ & $-2{,}105510$ & $-1{,}173280$ & $0{,}226718$ \\
        $4$ & $1{,}155150$ & $-2{,}183850$ & $-1{,}296890$ & $0{,}123602$ \\
        $5$ & $1{,}204800$ & $-2{,}243810$ & $-1{,}348720$ & $0{,}059968$ \\
        $6$ & $1{,}229590$ & $-2{,}284520$ & $-1{,}386680$ & $0{,}040706$ \\
        $7$ & $1{,}247280$ & $-2{,}304410$ & $-1{,}410490$ & $0{,}023809$ \\
        $8$ & $1{,}257540$ & $-2{,}318720$ & $-1{,}423350$ & $0{,}014307$ \\
        $9$ & $1{,}263610$ & $-2{,}327060$ & $-1{,}432050$ & $0{,}008709$ \\
        $10$ & $1{,}267520$ & $-2{,}331950$ & $-1{,}437150$ & $0{,}005095$ \\
        $11$ & $1{,}269810$ & $-2{,}335120$ & $-1{,}440210$ & $0{,}003171$ \\
        $12$ & $1{,}271220$ & $-2{,}336970$ & $-1{,}442150$ & $0{,}001933$ \\
        $13$ & $1{,}272090$ & $-2{,}338110$ & $-1{,}443290$ & $0{,}001146$ \\
        $14$ & $1{,}272610$ & $-2{,}338820$ & $-1{,}444000$ & $0{,}000705$ \\
        $15$ & $1{,}272930$ & $-2{,}339240$ & $-1{,}444430$ & $0{,}000432$ \\
        $16$ & $1{,}273130$ & $-2{,}339500$ & $-1{,}444690$ & $0{,}000260$ \\
        $17$ & $1{,}273250$ & $-2{,}339660$ & $-1{,}444850$ & $0{,}000160$ \\
        $18$ & $1{,}273320$ & $-2{,}339750$ & $-1{,}444950$ & $9{,}72 \cdot 10^{-5}$ \\
        $19$ & $1{,}273360$ & $-2{,}339810$ & $-1{,}445010$ & $5{,}92 \cdot 10^{-5}$ \\
        $20$ & $1{,}273390$ & $-2{,}339850$ & $-1{,}445040$ & $3{,}61 \cdot 10^{-5}$ \\
        $21$ & $1{,}273410$ & $-2{,}339870$ & $-1{,}445060$ & $2{,}20 \cdot 10^{-5}$ \\
        $22$ & $1{,}273420$ & $-2{,}339880$ & $-1{,}445080$ & $1{,}34 \cdot 10^{-5}$ \\
        $23$ & $1{,}273420$ & $-2{,}339890$ & $-1{,}445090$ & $8{,}17 \cdot 10^{-6}$ \\
        \hline
    \end{tabular}
\end{table}

На итерации №~23 достигнута разность $\|x^{(23)} - x^{(22)}\|_\infty = 8{,}17 \cdot 10^{-6} < 10^{-5}$, что завершило вычислительный процесс.

Итоговый вектор решения задачи 2 методом простой итерации:
\begin{equation}
    x_{\text{Iter}} = \begin{pmatrix} 1{,}27342 \\ -2{,}33989 \\ -1{,}44509 \end{pmatrix}.
    \label{eq:sol_iter}
\end{equation}
